\scanpage{23}
% archive sheet number: 22
On généralise la formule fondamentale au cas d'un \uline{morphisme} \uline{propre} \uline{quelconque}
\[
 f:X\longrightarrow Y
\]
\unsure{de} \deleted{\unsure{deux}} variétés \deleted{\unsure{analytiques|algébriques}} holomorphes, \unsure{partout|partant|purement} de codim relative $p$ :
\diagram{\fbox{\begin{tikzcd}[ampersand replacement=\&,column sep=huge,row sep=large]
 {H^{\deleted{\unsure{2i}}\unsure{i|*|2i}}(X,\mathbb C)}
 \arrow[r,"{f_*^{\mathrm{top}}}"] \arrow[d,dash,"{\wr}"'] \&
 {H^{\deleted{\unsure{2i}}\unsure{i+p|i+2p|2i-2p}}(Y,\mathbb C)} \arrow[d,dash,"{\wr}"']\\
 {H^{\deleted{\unsure{2}}i}(X,\Omega_X^\cdot)}
 \arrow[r,"{\frac{1}{(2i\pi)^p}f_*^{\mathrm{an}}}"] \&
 {H^{\deleted{\unsure{2}}i+\unsure{p|2p}}(Y,\Omega_Y^\cdot)}
\end{tikzcd}}}

D'où généralisation analogue pour des variétés \unsure{algébriques|algébriq.} sur $\mathbb C$ ....

